% Erklärung zur Konjunktion "bevor" und Vergleich mit der Präposition "vor"

\documentclass[a4paper,11pt]{article}

\usepackage[a4paper,margin=2cm]{geometry}
\usepackage[T1]{fontenc}
\usepackage[utf8]{inputenc}
\usepackage[ngerman,english]{babel}
\usepackage{lmodern}
\usepackage{microtype}
\usepackage[table]{xcolor}
\usepackage{parskip}
\usepackage{enumitem}
\usepackage{array}
\usepackage{longtable}
\usepackage[most]{tcolorbox}
\usepackage{titlesec}
\usepackage[hidelinks]{hyperref}

\definecolor{HeaderBlueDark}{RGB}{70,110,170}
\definecolor{RowBlueLight}{RGB}{235,243,255}
\definecolor{RowBlueDark}{RGB}{220,233,250}
\definecolor{SoftGray}{RGB}{90,90,90}

\setlength{\parindent}{0pt}
\setlist[itemize]{leftmargin=1.4em,itemsep=0.35em,topsep=0.35em}
\renewcommand{\arraystretch}{1.45}
\setlength{\tabcolsep}{6pt}
\linespread{1.06}

\titleformat{\section}[block]
{\Large\bfseries\color{HeaderBlueDark}}
{}{0pt}{}
\titlespacing{\section}{0pt}{1.3em}{0.8em}

\titleformat{\subsection}[block]
{\large\bfseries\color{HeaderBlueDark}}
{}{0pt}{}
\titlespacing{\subsection}{0pt}{1.0em}{0.6em}

\newtcolorbox{exbox}{enhanced,breakable,colback=RowBlueLight,colframe=RowBlueDark,boxrule=0.4pt,arc=1.5mm,left=1.2mm,right=1.2mm,top=1mm,bottom=1mm,title={Beispiele \textnormal{(Examples)}},fonttitle=\bfseries,colbacktitle=RowBlueDark,coltitle=black}

\NewDocumentEnvironment{grammarentry}{mm+b}{%
    \begin{tcolorbox}[enhanced,breakable,colback=white,colframe=HeaderBlueDark,boxrule=0.6pt,arc=1.5mm,left=1.5mm,right=1.5mm,top=1mm,bottom=1mm,colbacktitle=HeaderBlueDark,coltitle=white,fonttitle=\bfseries,title={#1\hfill\normalfont\footnotesize #2}]
    #3
    \end{tcolorbox}
}{}

\newcommand{\GermanText}[1]{\textbf{Deutsch: }#1\par}
\newcommand{\EnglishText}[1]{{\small\color{SoftGray}\textbf{English: }#1\par}}
\newcommand{\ExampleItem}[2]{\item \textbf{#1}\\{\small\color{SoftGray}#2}}

\title{Die Konjunktion \glqq bevor\grqq{} und der Vergleich mit \glqq vor\grqq}
\date{}

\begin{document}
    \maketitle
    \tableofcontents
    \newpage

    \section{Grundbedeutung von \glqq bevor\grqq}

        \begin{grammarentry}{bevor}{temporale, subordinierende Konjunktion}
            \GermanText{\textbf{bevor} bedeutet, dass eine Handlung früher als eine andere Handlung geschieht.}
            \EnglishText{\textbf{bevor} means that one action happens earlier than another action.}

            \GermanText{\textbf{bevor} verbindet einen Hauptsatz mit einem Nebensatz. Im Nebensatz steht das konjugierte Verb am Ende.}
            \EnglishText{\textbf{bevor} connects a main clause with a subordinate clause. In the subordinate clause, the conjugated verb stands at the end.}

            \GermanText{Grundstruktur: \textbf{Hauptsatz, bevor + Subjekt + weitere Satzteile + Verb.}}
            \EnglishText{Basic structure: \textbf{main clause, before + subject + other sentence elements + verb.}}
        \end{grammarentry}

        \begin{exbox}
            \begin{itemize}
                \ExampleItem{Ich wasche meine Hände, bevor ich esse.}{I wash my hands before I eat.}
                \ExampleItem{Sie prüft die Adresse, bevor sie den Brief abschickt.}{She checks the address before she sends the letter.}
                \ExampleItem{Wir kaufen die Fahrkarten, bevor wir in den Zug steigen.}{We buy the tickets before we get on the train.}
            \end{itemize}
        \end{exbox}

    \section{Zeitliche Beziehung der beiden Handlungen}

        \begin{grammarentry}{Welche Handlung geschieht zuerst?}{zeitliche Reihenfolge}
            \GermanText{Die Handlung im Hauptsatz geschieht normalerweise zuerst. Die Handlung im \textbf{bevor}-Satz geschieht danach.}
            \EnglishText{The action in the main clause normally happens first. The action in the \textbf{bevor} clause happens afterward.}

            \GermanText{Im Satz \glqq Ich wasche meine Hände, bevor ich esse\grqq{} geschieht das Händewaschen zuerst und das Essen danach.}
            \EnglishText{In the sentence ``I wash my hands before I eat,'' washing the hands happens first and eating happens afterward.}
        \end{grammarentry}

    \section{Stellung des \glqq bevor\grqq-Satzes}

        \begin{grammarentry}{Nebensatz nach dem Hauptsatz}{normale Reihenfolge}
            \GermanText{Wenn der \textbf{bevor}-Satz nach dem Hauptsatz steht, steht vor \textbf{bevor} ein Komma.}
            \EnglishText{When the \textbf{bevor} clause comes after the main clause, a comma stands before \textbf{bevor}.}
        \end{grammarentry}

        \begin{exbox}
            \begin{itemize}
                \ExampleItem{Ich rufe dich an, bevor ich losfahre.}{I call you before I leave.}
                \ExampleItem{Er liest den Vertrag, bevor er ihn unterschreibt.}{He reads the contract before he signs it.}
            \end{itemize}
        \end{exbox}

        \begin{grammarentry}{Nebensatz vor dem Hauptsatz}{Verb im Hauptsatz an zweiter Stelle}
            \GermanText{Der \textbf{bevor}-Satz kann auch am Anfang stehen. Danach folgt ein Komma. Im folgenden Hauptsatz steht das konjugierte Verb direkt nach dem Komma.}
            \EnglishText{The \textbf{bevor} clause can also stand at the beginning. A comma follows it. In the following main clause, the conjugated verb stands directly after the comma.}
        \end{grammarentry}

        \begin{exbox}
            \begin{itemize}
                \ExampleItem{Bevor ich losfahre, rufe ich dich an.}{Before I leave, I call you.}
                \ExampleItem{Bevor er den Vertrag unterschreibt, liest er ihn.}{Before he signs the contract, he reads it.}
            \end{itemize}
        \end{exbox}

    \section{Grundbedeutung von \glqq vor\grqq}

        \begin{grammarentry}{vor}{temporale Präposition mit Dativ}
            \GermanText{Bei einer Zeitangabe ist \textbf{vor} eine Präposition. Danach steht normalerweise ein Nomen oder eine Nomengruppe im Dativ.}
            \EnglishText{With a time expression, \textbf{vor} is a preposition. It is normally followed by a noun or noun phrase in the dative case.}

            \GermanText{Grundstruktur: \textbf{vor + Dativ}.}
            \EnglishText{Basic structure: \textbf{before + dative}.}

            \GermanText{\textbf{vor} leitet keinen Nebensatz ein. Deshalb steht direkt nach \textbf{vor} normalerweise kein vollständiger Satz mit Subjekt und konjugiertem Verb.}
            \EnglishText{\textbf{vor} does not introduce a subordinate clause. Therefore, it is normally not followed directly by a complete clause with a subject and a conjugated verb.}
        \end{grammarentry}

        \begin{exbox}
            \begin{itemize}
                \ExampleItem{Vor dem Essen wasche ich meine Hände.}{Before the meal, I wash my hands.}
                \ExampleItem{Sie prüft die Adresse vor dem Abschicken des Briefes.}{She checks the address before sending the letter.}
                \ExampleItem{Wir kaufen die Fahrkarten vor der Abfahrt.}{We buy the tickets before the departure.}
            \end{itemize}
        \end{exbox}

    \section{Vergleich zwischen \glqq bevor\grqq{} und \glqq vor\grqq}

        \rowcolors{2}{RowBlueLight}{RowBlueDark}
        \begin{longtable}{>{\bfseries}p{3.1cm}|p{5.2cm}|p{5.2cm}}
            \rowcolor{HeaderBlueDark}
            \color{white}\textbf{Merkmal} &
            \color{white}\textbf{bevor} &
            \color{white}\textbf{vor} \\
            Wortart & subordinierende Konjunktion & Präposition \\
            Englische Bedeutung & before & before \\
            Was folgt danach? & ein Nebensatz mit Subjekt und Verb & ein Nomen oder eine Nomengruppe \\
            Verbposition & das konjugierte Verb steht am Ende des Nebensatzes & keine besondere Verbposition; \textbf{vor} bildet keinen Nebensatz \\
            Kasus & kein Kasus & bei zeitlicher Bedeutung: Dativ \\
            Grundstruktur & bevor ich esse & vor dem Essen \\
            Beispielsatz & Ich wasche meine Hände, bevor ich esse. & Vor dem Essen wasche ich meine Hände. \\
        \end{longtable}

        \begin{grammarentry}{Hauptunterschied}{Satz oder Nomen}
            \GermanText{Benutze \textbf{bevor}, wenn danach ein vollständiger Nebensatz mit einem eigenen Subjekt und Verb folgt.}
            \EnglishText{Use \textbf{bevor} when it is followed by a complete subordinate clause with its own subject and verb.}

            \GermanText{Benutze \textbf{vor}, wenn danach ein Nomen oder eine Nomengruppe folgt.}
            \EnglishText{Use \textbf{vor} when it is followed by a noun or noun phrase.}
        \end{grammarentry}

    \section{Umformung von \glqq vor\grqq{} zu \glqq bevor\grqq}

        \begin{grammarentry}{Nomengruppe und Nebensatz}{gleiche zeitliche Bedeutung}
            \GermanText{Eine Verbindung mit \textbf{vor + Nomen} kann häufig in einen \textbf{bevor}-Satz umgeformt werden. Dabei muss aus dem Nomen eine Handlung mit Subjekt und Verb werden.}
            \EnglishText{A construction with \textbf{vor + noun} can often be changed into a \textbf{bevor} clause. In this process, the noun must become an action with a subject and a verb.}
        \end{grammarentry}

        \rowcolors{2}{RowBlueLight}{RowBlueDark}
        \begin{longtable}{p{6.5cm}|p{7cm}}
            \rowcolor{HeaderBlueDark}
            \color{white}\textbf{vor + Nomen} &
            \color{white}\textbf{bevor + Nebensatz} \\
            Vor dem Essen wasche ich meine Hände. & Bevor ich esse, wasche ich meine Hände. \\
            Vor der Abfahrt kontrollieren wir die Fahrkarten. & Bevor wir abfahren, kontrollieren wir die Fahrkarten. \\
            Vor seinem Anruf schrieb er eine E-Mail. & Bevor er anrief, schrieb er eine E-Mail. \\
            Vor der Prüfung wiederholt sie den Stoff. & Bevor sie die Prüfung schreibt, wiederholt sie den Stoff. \\
        \end{longtable}

    \newpage

    \section{bevor, bis und nachdem: zeitliche Unterschiede}
        \begin{grammarentry}{bevor und bis}{Reihenfolge versus Endpunkt}
            \GermanText{\textbf{bevor} zeigt, was früher geschieht. \textbf{bis} nennt den Endpunkt einer Handlung oder eines Zustands.}
            \EnglishText{\textbf{bevor} indicates what happens earlier. \textbf{bis} marks the end point of an action or state.}
        \end{grammarentry}
        \begin{exbox}
            \begin{itemize}
                \ExampleItem{Ich esse, bevor du kommst.}{I eat before you arrive.}
                \ExampleItem{Ich warte, bis du kommst.}{I wait until you arrive.}
            \end{itemize}
        \end{exbox}
        \begin{grammarentry}{bevor und nachdem}{before versus after}
            \GermanText{\textbf{bevor} bedeutet \glqq before\grqq{}, \textbf{nachdem} bedeutet \glqq after\grqq{}. Bei \textbf{nachdem} ist auch die zeitliche Beziehung der Zeitformen wichtig.}
            \EnglishText{\textbf{bevor} means before; \textbf{nachdem} means after. With \textbf{nachdem}, the relationship between tenses is also important.}
        \end{grammarentry}
        \begin{exbox}
            \begin{itemize}
                \ExampleItem{Ich frühstücke, bevor ich zur Arbeit gehe.}{I have breakfast before I go to work.}
                \ExampleItem{Nachdem ich gefrühstückt habe, gehe ich zur Arbeit.}{After I have had breakfast, I go to work.}
            \end{itemize}
        \end{exbox}

    \section{Zeitformen und Zukunft bei \glqq bevor\grqq}
        \begin{grammarentry}{Kein festes Tempus}{present, past and future}
            \GermanText{\textbf{bevor} verlangt keine feste Zeitform. Für zukünftige Ereignisse verwendet man im Nebensatz häufig \textbf{Präsens}; Futur~I ist nicht notwendig. Auch Präteritum und Perfekt sind möglich.}
            \EnglishText{\textbf{bevor} does not require a fixed tense. For future events, the subordinate clause often uses the present tense; future tense is unnecessary. Simple past and present perfect are also possible.}
        \end{grammarentry}
        \begin{exbox}
            \begin{itemize}
                \ExampleItem{Bevor ich morgen losfahre, werde ich die Reifen prüfen.}{Before I leave tomorrow, I will check the tyres.}
                \ExampleItem{Bevor ich losfuhr, prüfte ich die Reifen.}{Before I left, I checked the tyres.}
                \ExampleItem{Bevor ich losgefahren bin, habe ich die Reifen geprüft.}{Before I left, I checked the tyres.}
            \end{itemize}
        \end{exbox}

    \section{Erweiterungen und Varianten von \glqq bevor\grqq}
        \begin{grammarentry}{kurz, lange, noch bevor}{time nuances}
            \GermanText{\textbf{kurz bevor}, \textbf{lange bevor} und \textbf{noch bevor} präzisieren den zeitlichen Abstand. Das Verb steht im Nebensatz weiter am Ende.}
            \EnglishText{\textbf{kurz bevor} (shortly before), \textbf{lange bevor} (long before) and \textbf{noch bevor} (even before) specify the time interval. The verb still goes at the end.}
        \end{grammarentry}
        \begin{exbox}
            \begin{itemize}
                \ExampleItem{Kurz bevor der Zug abfuhr, kam sie an.}{She arrived shortly before the train departed.}
                \ExampleItem{Noch bevor ich antworten konnte, legte er auf.}{He hung up before I could even answer.}
            \end{itemize}
        \end{exbox}
        \begin{grammarentry}{ehe und bevor}{near-synonyms}
            \GermanText{Die Konjunktion \textbf{ehe} bedeutet oft dasselbe wie \textbf{bevor}, klingt jedoch häufig etwas gehobener. Das Nomen \textbf{die Ehe} bedeutet Ehe oder Heirat.}
            \EnglishText{The conjunction \textbf{ehe} often means the same as \textbf{bevor}, but sounds more formal. The noun \textbf{die Ehe} means marriage.}
        \end{grammarentry}
        \begin{exbox}
            \begin{itemize}
                \ExampleItem{Überlege gut, ehe du entscheidest.}{Think carefully before you decide.}
                \ExampleItem{Bevor ich es vergesse: Morgen habe ich einen Termin.}{Before I forget: I have an appointment tomorrow.}
            \end{itemize}
        \end{exbox}
        \begin{grammarentry}{nicht \ldots{} bevor}{not until}
            \GermanText{In einem verneinten Satz bedeutet \textbf{nicht \ldots{} bevor} oft \glqq erst wenn\grqq{}. Die Form \textbf{erst wenn} kann Missverständnisse vermeiden.}
            \EnglishText{In a negative sentence, not \ldots{} before often means not until. \textbf{erst wenn} can avoid ambiguity.}
        \end{grammarentry}
        \begin{exbox}
            \begin{itemize}
                \ExampleItem{Ich gehe nicht, bevor du zurück bist.}{I will not leave until you are back.}
                \ExampleItem{Ich gehe erst, wenn du zurück bist.}{I will leave only when you are back.}
            \end{itemize}
        \end{exbox}
    \section{Häufige Fehler}

        \begin{grammarentry}{Falsch: \glqq vor ich gehe\grqq}{nach \glqq vor\grqq{} kein Nebensatz}
            \GermanText{Falsch: \textbf{Vor ich gehe, rufe ich dich an.}}
            \EnglishText{Incorrect: \textbf{Before I go, I call you.}}

            \GermanText{Richtig mit Nebensatz: \textbf{Bevor ich gehe, rufe ich dich an.}}
            \EnglishText{Correct with a subordinate clause: \textbf{Before I go, I call you.}}

            \GermanText{Richtig mit Nomen: \textbf{Vor meinem Weggehen rufe ich dich an.}}
            \EnglishText{Correct with a noun: \textbf{Before my departure, I call you.}}
        \end{grammarentry}

        \begin{grammarentry}{Falsch: \glqq bevor dem Essen\grqq}{nach \glqq bevor\grqq{} braucht man einen Nebensatz}
            \GermanText{Falsch: \textbf{Bevor dem Essen wasche ich meine Hände.}}
            \EnglishText{Incorrect: \textbf{Before the meal, I wash my hands.}}

            \GermanText{Richtig mit Nomen: \textbf{Vor dem Essen wasche ich meine Hände.}}
            \EnglishText{Correct with a noun: \textbf{Before the meal, I wash my hands.}}

            \GermanText{Richtig mit Nebensatz: \textbf{Bevor ich esse, wasche ich meine Hände.}}
            \EnglishText{Correct with a subordinate clause: \textbf{Before I eat, I wash my hands.}}
        \end{grammarentry}

        \begin{grammarentry}{Falsche Verbposition}{Verb am Ende des Nebensatzes}
            \GermanText{Falsch: \textbf{Ich rufe dich an, bevor ich gehe nach Hause.}}
            \EnglishText{Incorrect: \textbf{I call you before I go home.}}

            \GermanText{Richtig: \textbf{Ich rufe dich an, bevor ich nach Hause gehe.}}
            \EnglishText{Correct: \textbf{I call you before I go home.}}
        \end{grammarentry}


    % Erweiterung C1-C2: fortgeschrittene Grammatik und Verwendung
    \section{Textorganisation und Satzstellung (Discourse organisation and word order)}
        \begin{grammarentry}{bevor als Nebensatzeinleiter}{subordinating conjunction}
            \GermanText{\textbf{bevor} leitet einen temporalen Adverbialsatz ein. Dieser kann dem Hauptsatz vorausgehen, folgen oder als Einschub stehen. Bei einer Voranstellung folgt das finite Hauptsatzverb direkt auf das Komma. Bei einer Einschiebung wird der Nebensatz durch zwei Kommas abgegrenzt.}
            \EnglishText{\textbf{bevor} introduces a temporal adverbial clause. It can precede, follow or interrupt the main clause. With initial placement, the finite main-clause verb follows the comma immediately. An inserted clause is enclosed in commas.}
        \end{grammarentry}
        \begin{exbox}
            \begin{itemize}
                \ExampleItem{Wir sollten, bevor wir endgültig entscheiden, weitere Gutachten einholen.}{Before reaching a final decision, we should obtain further expert opinions.}
                \ExampleItem{Bevor wir zur Abstimmung kommen, klären wir die offenen Fragen.}{Before we proceed to the vote, we will clarify the unresolved issues.}
            \end{itemize}
        \end{exbox}
        \begin{grammarentry}{Gesprächssteuerung}{discourse function}
            \GermanText{Bei \textbf{Bevor ich darauf eingehe, \ldots{}} und \textbf{Bevor ich es vergesse, \ldots{}} geht es nicht nur um Chronologie. Die Konstruktion kann einen Themenwechsel, eine Vorbemerkung oder die Reihenfolge argumentativer Schritte markieren.}
            \EnglishText{Expressions such as \textbf{Bevor ich darauf eingehe, \ldots{}} and \textbf{Bevor ich es vergesse, \ldots{}} do more than express chronology: they can signal a topic change, preliminary remark or sequence of argumentative steps.}
        \end{grammarentry}
        \begin{exbox}
            \begin{itemize}
                \ExampleItem{Bevor ich auf Ihren Einwand eingehe, möchte ich eine Voraussetzung klären.}{Before addressing your objection, I would like to clarify an assumption.}
                \ExampleItem{Bevor ich es vergesse: Die Frist endet am Freitag.}{Before I forget: the deadline ends on Friday.}
            \end{itemize}
        \end{exbox}

    \section{Ereignisfolge und Nicht-Eintreten (Sequence and non-occurrence)}
        \begin{grammarentry}{Keine notwendige Realisierung beider Ereignisse}{non-factual subordinate event}
            \GermanText{Ein \textbf{bevor}-Satz ordnet normalerweise zwei Ereignisse zeitlich. Trotzdem behauptet die Konstruktion \textbf{nicht immer}, dass das zweite Ereignis eingetreten ist: Eine Handlung kann verhindert werden, bevor sie überhaupt möglich wird.}
            \EnglishText{A \textbf{bevor}-clause normally orders two events in time, but it does \textbf{not always} imply that the later event actually occurred: an action may be prevented before it becomes possible.}
        \end{grammarentry}
        \begin{exbox}
            \begin{itemize}
                \ExampleItem{Der Antrag wurde zurückgezogen, bevor darüber abgestimmt werden konnte.}{The motion was withdrawn before it could be voted on.}
                \ExampleItem{Bevor er auch nur ein Wort sagen konnte, wurde die Sitzung beendet.}{Before he could even say a word, the meeting was ended.}
            \end{itemize}
        \end{exbox}
        \begin{grammarentry}{Zeitliche Relation bei Umstellung}{temporal interpretation}
            \GermanText{Die zeitliche Relation bleibt erhalten, auch wenn Haupt- und Nebensatz ihre Positionen tauschen. Sie richtet sich nicht nach der Reihenfolge der Satzteile.}
            \EnglishText{The temporal relation stays the same even when the main and subordinate clauses switch positions. It is not determined by their word order.}
        \end{grammarentry}
        \begin{exbox}
            \begin{itemize}
                \ExampleItem{Sie hatte das Gebäude verlassen, bevor der Alarm ausgelöst wurde.}{She had left the building before the alarm was triggered.}
                \ExampleItem{Bevor der Alarm ausgelöst wurde, hatte sie das Gebäude verlassen.}{Before the alarm was triggered, she had left the building.}
            \end{itemize}
        \end{exbox}

    \section{Tempus, Modus und Negation (Tense, mood and negation)}
        \begin{grammarentry}{Zeitformen im Kontext}{tense relationships}
            \GermanText{Die Zeitformen folgen dem zeitlichen Kontext, nicht einer starren Regel für \textbf{bevor}. In der Vergangenheit kann das Plusquamperfekt im Hauptsatz die bereits abgeschlossene Handlung markieren. Zukünftige Ereignisse können im Nebensatz im Präsens stehen. Der Konjunktiv II ist in hypothetischen Kontexten möglich.}
            \EnglishText{Tenses depend on context, not on a fixed tense rule for \textbf{bevor}. In past narration, pluperfect in the main clause may mark an already completed action. A future event can appear in the present tense in the subordinate clause. Subjunctive II is possible in hypothetical contexts.}
        \end{grammarentry}
        \begin{exbox}
            \begin{itemize}
                \ExampleItem{Bevor die Entscheidung bekannt gegeben wurde, waren alle Beteiligten informiert worden.}{Before the decision was announced, all those involved had been informed.}
                \ExampleItem{Ich informiere Sie, bevor die Änderungen in Kraft treten.}{I will inform you before the changes take effect.}
                \ExampleItem{Ich würde den Entwurf prüfen, bevor ich ihn freigäbe.}{I would examine the draft before approving it.}
            \end{itemize}
        \end{exbox}
        \begin{grammarentry}{Nicht bevor, erst wenn}{scope of negation}
            \GermanText{\textbf{Ich gehe nicht, bevor er kommt} verneint das Gehen vor seiner Ankunft. Häufig ist damit gemeint, dass man bis zu seiner Ankunft bleibt. Bei rechtlichen oder technischen Bedingungen können \textbf{erst wenn}, \textbf{nicht vor} oder \textbf{frühestens nach} eindeutiger sein.}
            \EnglishText{\textbf{Ich gehe nicht, bevor er kommt} denies leaving before he arrives. It often means that the speaker will stay until his arrival. For legal or technical conditions, \textbf{erst wenn}, \textbf{nicht vor} or \textbf{frühestens nach} may be more precise.}
        \end{grammarentry}
        \begin{exbox}
            \begin{itemize}
                \ExampleItem{Wir unterschreiben nicht, bevor sämtliche Fragen geklärt sind.}{We will not sign until all questions have been resolved.}
                \ExampleItem{Wir unterschreiben erst, wenn sämtliche Fragen geklärt sind.}{We will sign only once all questions have been resolved.}
            \end{itemize}
        \end{exbox}

    \section{Alternative Konstruktionen (Alternative constructions)}
        \begin{grammarentry}{ehe, bis, nachdem}{subtle contrasts}
            \GermanText{\textbf{ehe} ist oft eine gehobene oder literarisch wirkende Alternative zu \textbf{bevor}. \textbf{bis} kennzeichnet dagegen eine Grenze oder das Ende einer Dauer; \textbf{nachdem} kennzeichnet Nachzeitigkeit.}
            \EnglishText{\textbf{ehe} is often a more formal or literary alternative to \textbf{bevor}. \textbf{bis} instead marks a boundary or the end of a duration; \textbf{nachdem} marks posteriority.}
        \end{grammarentry}
        \begin{exbox}
            \begin{itemize}
                \ExampleItem{Noch ehe die Verhandlungen begannen, war eine Einigung ausgeschlossen.}{Even before the negotiations began, an agreement was ruled out.}
                \ExampleItem{Wir warteten, bis die Verhandlungen begannen.}{We waited until negotiations started.}
                \ExampleItem{Nachdem die Verhandlungen begonnen hatten, änderte sich die Lage.}{After negotiations had begun, the situation changed.}
            \end{itemize}
        \end{exbox}
        \begin{grammarentry}{lieber, eher als, bevor}{preference instead of mere chronology}
            \GermanText{In Kontexten der Wahl oder Ablehnung kann \textbf{bevor} auch eine Präferenz nahelegen: \textbf{Bevor ich nachgebe, gehe ich lieber.} Ein ausdrücklicher Vergleich lässt sich mit \textbf{lieber \ldots{} als} oder \textbf{eher \ldots{} als dass} formulieren.}
            \EnglishText{In choice or refusal contexts, \textbf{bevor} may imply a preference: rather than give in, the speaker would leave. Explicit comparison can also use \textbf{lieber \ldots{} als} or \textbf{eher \ldots{} als dass}.}
        \end{grammarentry}
        \begin{exbox}
            \begin{itemize}
                \ExampleItem{Bevor ich eine unbelegte Behauptung wiederhole, schweige ich lieber.}{Rather than repeat an unsupported claim, I would prefer to remain silent.}
                \ExampleItem{Ich schweige lieber, als eine unbelegte Behauptung zu wiederholen.}{I would rather remain silent than repeat an unsupported claim.}
            \end{itemize}
        \end{exbox}

    \section{Nominalisierung im formellen Stil (Nominalisation in formal style)}
        \begin{grammarentry}{bevor oder vor plus Nomen}{formal paraphrase}
            \GermanText{In Verwaltung, Wissenschaft und Recht lässt sich ein \textbf{bevor}-Satz oft durch eine Nominalgruppe mit \textbf{vor + Dativ} ersetzen. Der Nominalstil kann die handelnde Person in den Hintergrund treten lassen.}
            \EnglishText{In administration, academic prose and legal texts, a \textbf{bevor}-clause can often be replaced by a nominal phrase with \textbf{vor + dative}. Nominal style may place less emphasis on the acting person.}
        \end{grammarentry}
        \begin{exbox}
            \begin{itemize}
                \ExampleItem{Bevor die Verordnung in Kraft tritt, müssen die Behörden informiert werden.}{Before the regulation takes effect, the authorities must be informed.}
                \ExampleItem{Vor Inkrafttreten der Verordnung müssen die Behörden informiert werden.}{The authorities must be informed before the regulation comes into force.}
                \ExampleItem{Bevor der Vertrag unterzeichnet wird, sind die Bedingungen zu prüfen.}{Before the contract is signed, the terms must be reviewed.}
                \ExampleItem{Vor Unterzeichnung des Vertrags sind die Bedingungen zu prüfen.}{Before the contract is signed, the terms must be reviewed.}
            \end{itemize}
        \end{exbox}
    \section{Zusammenfassung}

        \begin{grammarentry}{Merksatz}{kurz und wichtig}
            \GermanText{\textbf{bevor + Satz}: Bevor ich esse, wasche ich meine Hände.}
            \EnglishText{\textbf{before + clause}: Before I eat, I wash my hands.}

            \GermanText{\textbf{vor + Nomen im Dativ}: Vor dem Essen wasche ich meine Hände.}
            \EnglishText{\textbf{before + noun in the dative}: Before the meal, I wash my hands.}

            \GermanText{Merke: \textbf{bevor} braucht einen Nebensatz; temporales \textbf{vor} braucht normalerweise ein Nomen im Dativ.}
            \EnglishText{Remember: \textbf{bevor} requires a subordinate clause; temporal \textbf{vor} normally requires a noun in the dative.}
        \end{grammarentry}

\end{document}
